(1) is Theorem \claimref{singularity.kms-uniqueness-factor}, read over the atom's scalars $\mathrm{span}\{1, J\}$ (its unit $I_4$ there standing for the factor's scalar unit, here $J$), with $K$ replaced by $\omega K$, and $Z = \tr e^{-\beta\omega K} = e^{\beta\omega/2}(1 + e^{-\beta\omega})$, the prefactor cancelling in the normalised state. (2) Write $\phi = \tr(\varrho\,\cdot\,)$ and $H := \omega K$. The graded identity with $b$ homogeneous is $\tr(b\varrho a) = \epsilon_b\tr(e^{\beta H}\varrho\,b\,e^{-\beta H}a)$ for all $a$, i.e.\ $b\,(\varrho e^{\beta H}) = \epsilon_b\,(e^{\beta H}\varrho)\,b$. With $b = 1$: $\varrho$ commutes with $e^{\beta H}$, hence with $H$, hence is diagonal. Then $c := \varrho e^{\beta H}$ is diagonal, commutes with the even elements and anticommutes with $v$ and $v^{*}$; for $c = \operatorname{diag}(c_0, c_1)$ the relation $cv = -vc$ reads $c_0 = -c_1$, so $c \propto \Gamma$ and $\varrho \propto \Gamma e^{-\beta H}$. Conversely $\Gamma e^{-\beta H}$ satisfies the graded identity because $\Gamma$ commutes with $H$ and $b\Gamma = \epsilon_b\Gamma b$. Its trace is $e^{\beta\omega/2}(1 - e^{-\beta\omega})$ and $\tr(\Gamma e^{-\beta H}N) = -e^{-\beta\omega}e^{\beta\omega/2}$, giving the displayed ratio. (3) is read off (1) and (2): $\tr(e^{-\beta H}) $ and $\tr(\Gamma e^{-\beta H})$ differ exactly in the sign of the occupied term. The script verifies both identities on all sixteen matrix-unit pairs, the failure of each functional to satisfy the other identity, and the negativity of the twisted occupancy, at $e^{-\beta\omega} = \tfrac{9}{16}$.
